\documentclass[10pt,a4paper]{article}
\usepackage[utf8]{inputenc}
\usepackage[a4paper,left=3cm,right=2cm]{geometry}
\usepackage{amsmath}
\usepackage[thmmarks, amsmath, thref]{ntheorem}
\usepackage{amsfonts}
\usepackage{amssymb}
\usepackage{fancyhdr}
\pagestyle{fancy}
\usepackage{anysize}
\usepackage{graphicx}
\usepackage{float}
\usepackage{hyperref}
\usepackage{multirow}
\usepackage{enumitem}
\newtheorem{theorem}{Theorem}
\theoremsymbol{\ensuremath{\blacksquare}}
\newtheorem*{solution}{Solución:}

\usepackage{xcolor}

% Margenes y formalidades

\textwidth 6.25in % Margen Derecho y ancho del texto
\textheight 8.5in
\oddsidemargin 0in % Margen Izquierdo
\headheight 0in % Largo del texto
\leftmargin 1in
\parindent 0em % Salto entre la indentación (Sangría)https://www.overleaf.com/4519127669sdpptkhvsbfn
\topmargin -0.5in % Margen Superior
\parskip 0.5ex % Salto entre parrafos
\baselineskip 1pt

%\lhead{\includegraphics[scale=0.2]{logo.png}} % texto izquierda de la cabecera
%\chead{}                                      % texto centro de la cabecera
%\rhead{\includegraphics[scale=0.2]{FIG/logo-mba-2016.png}} % número de página a la derecha

\renewcommand{\headrulewidth}{0.4pt} % grosor de la línea de la cabecera
\renewcommand{\footrulewidth}{0.4pt} % grosor de la línea del pie

\begin{document}



\pagebreak
\vspace*{4mm} 

\begin{center}
\begin{huge}
\textbf{\\Ejercicios Econometría 2}
\end{Large}\\
%{\large 2023-2}\\
\end{center}


\section{Muestreo e Intervalos de Confianza}

Un investigador quiere estimar el gasto promedio mensual en transporte de los estudiantes de una universidad. Para ello, selecciona una muestra aleatoria de 100 estudiantes y obtiene los siguientes datos:

\begin{itemize}
    \item \textbf{Media poblacional} ($\mu$) = 120 dólares (dado como referencia)
    \item \textbf{Varianza poblacional} ($\sigma^2$) = 900 dólares
    \item \textbf{Media muestral} ($\bar{x}$) = 125 dólares
    \item \textbf{Tamaño de la muestra} ($n$) = 100
    \item \textbf{Nivel de confianza} = 95\%
\end{itemize}

\textbf{Preguntas:}

\begin{enumerate}
    \item Calcula el intervalo de confianza del 95\% para la media poblacional.
    \item Interpreta el intervalo de confianza en el contexto del problema.
    \item Determina el tamaño mínimo de la muestra ($n$) necesario para que la media poblacional de 120 no pertenezca al intervalo de confianza del 95\% basado en la media muestral de 125.
\end{enumerate}

\textbf{Solución}

\textbf{Intervalo de confianza}
El intervalo de confianza para la media poblacional se calcula con la fórmula:

\begin{equation}
IC = \left( \bar{x} \pm Z_{\alpha/2} \cdot \frac{\sigma}{\sqrt{n}} \right)
\end{equation}

\begin{itemize}
    \item \textbf{Desviación estándar poblacional:} 
    \begin{equation}
    \sigma = \sqrt{\sigma^2} = \sqrt{900} = 30
    \end{equation}
    \item \textbf{Error estándar de la media:}
    \begin{equation}
    SE = \frac{\sigma}{\sqrt{n}} = \frac{30}{\sqrt{100}} = \frac{30}{10} = 3
    \end{equation}
    \item \textbf{Valor crítico para un 95\% de confianza:}
    \begin{equation}
    Z_{\alpha/2} = 1.96
    \end{equation}
    \item \textbf{Cálculo del intervalo:}
    \begin{equation}
    IC = \left( 125 \pm 1.96 \cdot 3 \right)
    \end{equation}
    \begin{equation}
    IC = \left( 125 \pm 5.88 \right)
    \end{equation}
    \begin{equation}
    IC = (119.12, 130.88)
    \end{equation}
\end{itemize}

\textbf{Interpretación}
Con un 95\% de confianza, el verdadero gasto promedio mensual en transporte de los estudiantes está entre \textbf{119.12 y 130.88 dólares}. Esto significa que si repitiéramos este muestreo muchas veces, el 95\% de los intervalos construidos contendrían la verdadera media poblacional.\\

\textbf{Tamaño mínimo de la muestra ($n$)}
Queremos encontrar el mínimo $n$ tal que $\mu = 120$ no pertenezca al intervalo de confianza. Es decir, buscamos que:

\begin{equation}
\bar{x} - Z_{\alpha/2} \cdot \frac{\sigma}{\sqrt{n}} > 120
\end{equation}

Sustituyendo los valores conocidos:

\begin{equation}
125 - 1.96 \cdot \frac{30}{\sqrt{n}} > 120
\end{equation}

Restamos 125 en ambos lados:

\begin{equation}
- 1.96 \cdot \frac{30}{\sqrt{n}} > -5
\end{equation}

Multiplicamos por $-1$ (cambiando el sentido de la desigualdad):

\begin{equation}
1.96 \cdot \frac{30}{\sqrt{n}} < 5
\end{equation}

Dividimos por 1.96:

\begin{equation}
\frac{30}{\sqrt{n}} < \frac{5}{1.96}
\end{equation}

Aproximamos:

\begin{equation}
\frac{30}{\sqrt{n}} < 2.55
\end{equation}

Invirtiendo la ecuación:

\begin{equation}
\sqrt{n} > \frac{30}{2.55}
\end{equation}

\begin{equation}
\sqrt{n} > 11.76
\end{equation}

Elevamos al cuadrado:

\begin{equation}
n > 138.3
\end{equation}

Como $n$ debe ser un número entero, tomamos $n = 139$.

\section{Regresión Lineal Simple con suma de Constante}

Sean \(\hat{\beta}_0\) y \(\hat{\beta}_1\) el intercepto y la pendiente en la regresión de \(y_i\) sobre \(x_i\), usando \(n\) observaciones.

Ahora, sean \(\tilde{\beta}_0\) y \(\tilde{\beta}_1\) la intersección con el eje \(y\) y la pendiente en la regresión de \((c_1 + y_i)\) sobre \((c_2 + x_i)\) (sin ninguna restricción sobre \(c_1\) o \(c_2\)). Muestre que \(\tilde{\beta}_1 = \hat{\beta}_1\) y \(\tilde{\beta}_0 = \hat{\beta}_0 +c_1- c_2 \cdot \hat{\beta}_1\). Compare los $R^2$ de ambas regresiones.

\textbf{Solución}
{ Estimadores MCO.

$$\hat{\beta}_0 = \bar{y} - \hat{\beta}_1 \bar{x} $$
$$\hat{\beta}_1 = \frac{\sum_{i=1}^{n}(x_i - \bar{x})(y_i-\bar{y})}{\sum_{i=1}^{n}(x_i - \bar{x})^2}$$


$$\tilde{\beta}_1 = \frac{\sum_{i=1}^{n}((x_i+c_2) - (\bar{x}+c_2))((y_i+c_1)-(\bar{y}+c_1))}{\sum_{i=1}^{n}((x_i+c_2) - (\bar{x}+c_2))^2}=\hat{\beta}_1$$

$$\tilde{\beta}_0 = (\bar{y}+c_1) - \tilde{\beta}_1 (\bar{x}+c_2)= (\bar{y}+c_1) - \hat{\beta}_1 (\bar{x}+c_2)=\hat{\beta}_0+c_1-c_2\hat{\beta}_1  $$
}



    
{ SST la suma de cuadrados totales (total sum of squares):
		\[ SST = \sum_{i=1}^{n} (y_i - \bar{y})^2 \]
		
	 SSE la suma de cuadrados explicados (explained sum of squares):
	{	\[ SSE = \sum_{i=1}^{n} (\hat{y}_i - \bar{y})^2 \]}	

 Para el segundo caso
		\[ SST' = \sum_{i=1}^{n} ((y_i+c_1) - (\bar{y}+c_1))^2 =SST\]
		
		
$$\tilde{y}_i=\tilde{\beta_0}+\tilde{\beta_1}(x_i+c_2)=\hat{\beta}_0+c_1-c_2\hat{\beta}_1 + \hat{\beta_1}(x_i+c_2)= \hat{\beta_0}+c_1+\hat{\beta_1}x_i$$
	{	\[ SSE' = \sum_{i=1}^{n} (\tilde{y}_i - (\bar{y}+c_1))^2=SSE \]

 $$SSE' \sum_{i=1}^{n} (\hat{\beta_0}+c_1+\hat{\beta_1}x_i - (\bar{y}+c_1))^2=SSE $$
 
 Luego el $R^2$ de ambos modelos es igual. 
 }		
 
 
 }



\section{Regresión Lineal Simple: Sesgo por error de medición}

Sean \( \hat{\beta}_0 \) y \( \hat{\beta}_1 \) los parámetros estimados de la regresión de \( y_i \) sobre \( x_i \), y sean \( \tilde{\beta}_0 \) y \( \tilde{\beta}_1 \) los parámetros estimados de la regresión de \( y_i \) sobre \( (u + x_i) \), donde \( u \sim N(0, 1) \). 
\begin{enumerate}
    \item Encuentre los nuevos valores de \( \tilde{\beta}_0 \) y \( \tilde{\beta}_1 \) en función de \( \hat{\beta}_0 \) y \( \hat{\beta}_1 \).

    \item Supongamos que tenemos una muestra de 1000 datos con $\Bar{x} = 5$, $\Bar{y}=10$, $Var(x) = 4$ y $Cov(x,y)=3$, encuentre el valor de los estimadores \( \hat{\beta}_0 \) y \( \hat{\beta}_1 \) y de \( \tilde{\beta}_0 \) y \( \tilde{\beta}_1 \), y las ecuaciones que describen al regresando.

\end{enumerate}

\textbf{Solución}
\begin{enumerate}
    \item \textbf{Modelo original}
El modelo original es:
\[
y_i = \beta_0 + \beta_1 x_i + \epsilon_i
\]
donde \( \epsilon_i \) es el error aleatorio. \\

\textbf{Nuevo modelo}
En el nuevo modelo, regresamos \( y_i \) sobre \( (u_i + x_i) \), lo que da la ecuación:
\[
y_i = \tilde{\beta}_0 + \tilde{\beta}_1 (u_i + x_i) + \nu_i
\]
donde \( u \sim N(0, 1) \) es una variable aleatoria con media cero y varianza uno. Usemos  $\tilde{x}_i = x_i + u_i$ \\

\textbf{Relación entre \( \hat{\beta}_0, \hat{\beta}_1 \) y \( \tilde{\beta}_0, \tilde{\beta}_1 \)}
Sabemos que:
\[
\hat{\beta}_1 = \frac{Cov(x,y)}{Var(x)} 
\]
y
\[
\hat{\beta}_0 = \bar{y} - \hat{\beta}_1 \bar{x} 
\]

para los nuevos valores tenemos que 
$$\tilde{\beta}_1 = \frac{Cov(\tilde{x},y)}{Var(\tilde{x})} = \frac{Cov(x+u,y)}{Var(x+u)}$$
Sabemos que $Cov(u,y)=0$, por lo que:
$$\tilde{\beta}_1 = \frac{Cov(x,y)}{Var(x+u)}$$
sustituyendo $y = \hat{\beta_0} + \hat{\beta_1}x +\epsilon $ tenemos:
$$\tilde{\beta}_1 = \frac{Cov(x,\hat{\beta_0} + \hat{\beta_1}x +\epsilon)}{Var(x+u)} = \hat{\beta_1} \frac{Cov(x,x)}{Var(x+u)} = \hat{\beta_1}\frac{Var(x)}{Var(x)+Var(u)}$$

para el valor de $\tilde{\beta}_0$, tenemos:
$$\tilde{\beta}_0 = \Bar{y} - \tilde{\beta}_1 \Bar{x}$$
reemplazando el valor de $\tilde{\beta}_1$, tenemos:
$$\tilde{\beta}_0 = \Bar{y} - \hat{\beta_1}\frac{Var(x)}{Var(x)+Var(u)} \Bar{x}$$
sabemos que $Var(u)=1$ entonces:
$$\tilde{\beta}_1 =  \hat{\beta_1}\frac{Var(x)}{Var(x)+1}$$
$$\tilde{\beta}_0 = \Bar{y} - \hat{\beta_1}\frac{Var(x)}{Var(x)+1} \Bar{x}$$

\item Usando los valores del enuncado y estos resultados tenemos que:
$$\hat{\beta}_1 = \frac{Cov(x,y)}{Var(x)}  = \frac{3}{4} = 0.75$$
$$\hat{\beta}_0 = \bar{y} - \hat{\beta}_1 \bar{x} = 10-0.75*5 = 6.25 $$
quedando la la ecuación del modelo original como:
$$y_i = 6.25 + 0.75 x_i + \epsilon_i$$
y reemplazando en los nuevos valores de los estimadores, tenemos:
$$\tilde{\beta}_1 =  \hat{\beta_1}\frac{Var(x)}{Var(x)+1} = 0.75*\frac{4}{4+1}=0.6$$
y como la media de $\tilde{x}$ es igual a la media de $x$, tenemos:
$$\tilde{\beta}_0 = \Bar{y} - \tilde{\beta_1} \Bar{x} = 10-0.6*5= 7$$
quedando la ecuación del modelo modificado como:
$$y_i = 7 + 0.6 (u_i + x_i) + \nu_i$$

\end{enumerate}


\section{Sesgo de selección}
Supongamos que un gobierno implementa un programa de capacitación laboral para ayudar a las personas desempleadas a encontrar trabajo. El programa ofrece cursos gratuitos de habilidades técnicas y entrevistas simuladas. Sin embargo, la participación en el programa es voluntaria, y las personas deciden inscribirse o no en función de sus propias características (por ejemplo, motivación, nivel educativo previo, etc.).

Tenemos los siguientes datos para dos grupos: \textbf{tratados} (participantes del programa) y \textbf{no tratados} (no participantes).

\begin{center}
\begin{tabular}{|c|c|c|c|}
\hline
Individuo & Tratamiento (\( D_i \)) & Salario (\( Y_i \)) & Nivel educativo (\( X_i \)) \\
\hline
1 & 1 & 2500 & 12 \\
2 & 1 & 3000 & 14 \\
3 & 1 & 2800 & 13 \\
4 & 0 & 2000 & 10 \\
5 & 0 & 2200 & 11 \\
6 & 0 & 2100 & 10 \\
\hline
\end{tabular}
\end{center}

\begin{enumerate}
    \item Comente cuales serían los posibles sesgos que existen en comparar el grupo de tratados vs los no tratados.

    \item Calcule el $\beta$ que asimila el ATE y ATT. 

    \item ¿Existe un sesgo dado el nivel educativo?

    \item supongamos que incluimos la variable nivel educativo en una regresión lineal dandonos la siguiente ecuación: 
    $$\hat{Y}_i = 1500+200D_i + 100X_i$$
    ¿Cúal es el valor del sesgo de selección? ¿Cúal sería el impacto real de la participación en el programa en el salario?
\end{enumerate}

\textbf{Solución}
\begin{enumerate}
    \item Queremos evaluar el efecto del programa en el salario de los participantes. Sin embargo, existe un \textbf{sesgo de selección} porque las personas que eligen participar en el programa pueden ser más motivadas o tener mejores habilidades que las que no participan, lo que podría afectar los resultados incluso si el programa no tuviera ningún efecto real.

    \item El efecto del tratamiento (\( \beta \)) es la diferencia en el salario promedio entre el grupo tratado y el grupo no tratado:
\[
\beta = \bar{Y}_{\text{tratado}} - \bar{Y}_{\text{no tratado}}.
\]

Calculamos los promedios:
\[
\bar{Y}_{\text{tratado}} = \frac{2500 + 3000 + 2800}{3} = \frac{8300}{3} = 2766.67,
\]
\[
\bar{Y}_{\text{no tratado}} = \frac{2000 + 2200 + 2100}{3} = \frac{6300}{3} = 2100.
\]

Por lo tanto:
\[
\beta = 2766.67 - 2100 = 666.67.
\]

Según este cálculo, el programa parece aumentar el salario en \textbf{\$666.67} en promedio.

\item Ahora, agregamos la información del nivel educativo (\( X_i \)) para ver si los grupos son comparables.

Calculamos las medias de los años de educación:
\[
\bar{X}_{\text{tratado}} = \frac{12 + 14 + 13}{3} = 13,
\]
\[
\bar{X}_{\text{no tratado}} = \frac{10 + 11 + 10}{3} = 10.33.
\]

\item Al ver la regresión
$$\hat{Y}_i = 1500+200D_i + 100X_i$$
Podemos ver que el $\beta_1=200$, por lo que el efecto que produce tomar o no el programa en el salario es de $200$. Usando el resultado del beta antes encontrado $\beta = 666.67$, podemos encontrar que el valor del sesgo de selección, en este caso es de $466.67$
$$666.67 = 200 + sesgo \implies sesgo = 466.67$$
\end{enumerate}



\section{Sesgo por variables omitidas}
Queremos estudiar el efecto de la educación (\( X_1 \)) sobre el salario (\( Y \)). Sin embargo, omitimos la variable "habilidades innatas" (\( X_2 \)), que también afecta el salario y está correlacionada con la educación. Esto introduce un \textbf{sesgo por variables omitidas}.

\textbf{Datos}
Tenemos los siguientes datos para 6 individuos:

\begin{center}
\begin{tabular}{|c|c|c|c|}
\hline
Individuo & Educación (\( X_1 \)) & Habilidades innatas (\( X_2 \)) & Salario (\( Y \)) \\
\hline
1 & 12 & 5 & 2500 \\
2 & 14 & 7 & 3000 \\
3 & 16 & 6 & 3500 \\
4 & 10 & 4 & 2000 \\
5 & 13 & 5 & 2800 \\
6 & 15 & 8 & 4000 \\
\hline
\end{tabular}
\end{center}

\begin{enumerate}
    \item Estimando el efecto de la educación (\( X_1 \)) sobre el salario (\( Y \)) omitiendo la variable "habilidades innatas" (\( X_2 \)) (regresión lineal simple), tenemos los siguientes resultados:
    \begin{itemize}
    \item Coeficiente de educación (\( \hat{\beta}_1 \)): 200.
    \item Intercepto (\( \hat{\beta}_0 \)): 500.
    \item \( R^2 \): 0.75.
    \item \( SSR \): 150,000.
    \end{itemize}
    Luego, incluyendo la variable "habilidades innatas" (\( X_2 \)) en el modelo (regresión múltiple), obtenemos los siguientes resultados:
    \begin{itemize}
    \item Coeficiente de educación (\( \hat{\beta}_1 \)): 150.
    \item Coeficiente de habilidades innatas (\( \hat{\beta}_2 \)): 100.
    \item Intercepto (\( \hat{\beta}_0 \)): 300.
    \item \( R^2 \): 0.90.
    \item \( SSR \): 50,000.
\end{itemize}
Compare los resultados de ambos modelos e identifique el signo del sesgo.

    \item Explica por qué ocurre el sesgo por variables omitidas.
\end{enumerate}

\textbf{Solución}

\begin{enumerate}
    \item Centralizando la información en una tabla tenemos:
    \begin{center}
\begin{tabular}{|l|c|c|}
\hline
Métrica & Regresión simple (sin \( X_2 \)) & Regresión múltiple (con \( X_2 \)) \\
\hline
Coeficiente de \( X_1 \) & 200 & 150 \\
Coeficiente de \( X_2 \) & No aplica & 100 \\
Intercepto (\( \beta_0 \)) & 500 & 300 \\
\( R^2 \) & 0.75 & 0.90 \\
\( SSR \) & 150,000 & 50,000 \\
\hline
\end{tabular}
\end{center}
\begin{itemize}
    \item El coeficiente de educación (\( X_1 \)) en la regresión simple (200) está \textbf{sobreestimado} en comparación con el de la regresión múltiple (150). Esto se debe a que la variable "habilidades innatas" (\( X_2 \)) no se incluyó en el primer modelo, y \( X_2 \) está correlacionada con \( X_1 \) y afecta a \( Y \).
    \item El \( R^2 \) aumenta de 0.75 a 0.90 al incluir \( X_2 \), lo que indica que el modelo con \( X_2 \) explica mejor la variabilidad del salario.
    \item El \( SSR \) disminuye de 150,000 a 50,000 al incluir \( X_2 \), lo que confirma un mejor ajuste del modelo.\\
    \\
    El signo del sesgo lo encontramos desde el $\beta_1$ de la regresión lineal simple que denotaremos $\tilde{\beta}_1$ y que se relaciona con los regreseros de la ecuación real (regresión lineal multiple) como:
    $$\tilde{\beta}_1 = \beta_1 + \beta_2 \delta_1$$
    Donde $\delta_1$ muestra la relación entre $X_1$ y $X_2$, entonces calculando tenemos:
    $$200 = 150 + 100\delta_1 \implies \delta_1 = 0.5 >0$$
    por lo que el sesgo es positivo.

    \item El sesgo por variables omitidas ocurre cuando excluimos una variable relevante (\( X_2 \)) que está correlacionada con la variable de interés (\( X_1 \)) y afecta a la variable dependiente (\( Y \)). En este ejercicio:
\begin{itemize}
    \item El coeficiente de educación (\( X_1 \)) está sobreestimado en la regresión simple debido a la omisión de "habilidades innatas" (\( X_2 \)).
    \item Al incluir \( X_2 \), el modelo mejora (\( R^2 \) aumenta y \( SSR \) disminuye), y el coeficiente de educación se ajusta a un valor más realista.
\end{itemize}
    
\end{itemize}


    
\end{enumerate}



\end{document}